\documentclass[11pt,a4paper]{article}
\usepackage[utf8]{inputenc}
\usepackage[vietnamese]{babel}
\usepackage{amsmath,amssymb}
\usepackage{geometry}
\geometry{a4paper, total={180mm,260mm}, left=15mm, right=15mm, top=18mm, bottom=18mm}

\usepackage{tikz}
\usetikzlibrary{shapes.geometric, arrows.meta, positioning, calc}
\tikzset{
  arr/.style={-{Stealth[scale=0.8]}},
  every node/.append style={align=center}
}
\usepackage{tcolorbox}
\tcbuselibrary{skins,breakable}
\usepackage{fontawesome}
\usepackage{fancyhdr}
\usepackage{enumitem}

\pagestyle{fancy}
\fancyhf{}
\renewcommand{\headrulewidth}{0.5pt}
\renewcommand{\footrulewidth}{0.5pt}
\fancyhead[L]{\small\bfseries\color{violet!85!black}\faBook\ SỔ TAY KIẾN THỨC TOÁN 2}
\fancyhead[R]{\small\bfseries\color{violet!85!black}Lớp Toán Thầy Lĩnh -- loptoanthaylinh.vn}
\fancyfoot[C]{\small\color{gray}-- Sổ tay kiến thức: Bài 40. Bảng nhân 5 --}

\setlength{\parindent}{0pt}
\setlength{\parskip}{0pt}

\begin{document}

% --- BANNER BÀI HỌC ---
\begin{tcolorbox}[
  enhanced,
  colback=violet!8!white,
  colframe=violet!70!black,
  arc=3mm, boxrule=1.2pt,
  left=3mm, right=3mm, top=2mm, bottom=2mm,
  halign=center
]
  {\LARGE\bfseries\color{violet!90!black}\faBookmark\ BÀI 40: BẢNG NHÂN 5}\\[1mm]
  {\itshape\footnotesize\color{black!75} Hình thành và ghi nhớ bảng cửu chương nhân 5; quy luật đếm thêm 5 trên 1 trang A4}
\end{tcolorbox}

\vspace{-2mm}

\noindent
\begin{minipage}[t]{0.49\textwidth}\vspace{0pt}
\begin{tcolorbox}[
  enhanced,
  colback=violet!4!white,
  colframe=violet!75!black,
  arc=2mm, boxrule=0.8pt,
  left=2mm, right=2mm, top=1.8mm, bottom=1.8mm,
  title={\footnotesize\bfseries\color{white}\faKey\ 1. Bảng nhân 5 chuẩn mực}
]
{\footnotesize
\begin{tabular}{l@{\quad}l}
  $5 \times 1 = \mathbf{5}$ & $5 \times 6 = \mathbf{30}$ \\
  $5 \times 2 = \mathbf{10}$ & $5 \times 7 = \mathbf{35}$ \\
  $5 \times 3 = \mathbf{15}$ & $5 \times 8 = \mathbf{40}$ \\
  $5 \times 4 = \mathbf{20}$ & $5 \times 9 = \mathbf{45}$ \\
  $5 \times 5 = \mathbf{25}$ & $5 \times 10 = \mathbf{50}$
\end{tabular}
\vspace{1mm}
\begin{itemize}\setlength{\itemsep}{0.5pt}\setlength{\leftskip}{-2.5mm}
  \item Thừa số thứ nhất đều là $5$.
  \item \textbf{Đặc điểm dãy tích:} Các số chỉ tận cùng bằng chữ số \textbf{0} hoặc \textbf{5}.
\end{itemize}
}
\end{tcolorbox}
\end{minipage}
\hfill
\begin{minipage}[t]{0.49\textwidth}\vspace{0pt}
\begin{tcolorbox}[
  enhanced,
  colback=teal!4!white,
  colframe=teal!75!black,
  arc=2mm, boxrule=0.8pt,
  left=2mm, right=2mm, top=1.8mm, bottom=1.8mm,
  title={\footnotesize\bfseries\color{white}\faSliders\ 2. Quy luật đếm thêm 5}
]
{\footnotesize
\begin{itemize}\setlength{\itemsep}{1pt}\setlength{\leftskip}{-2.5mm}
  \item Trong bảng nhân 5, tích liền sau luôn \textbf{hơn tích liền trước 5 đơn vị}.
  $$\text{Tích sau} = \text{Tích trước} + 5$$
  \item \textbf{Liên hệ thực tế:}\\
  - Mỗi bàn tay có $5$ ngón tay.\\
  - Khoảng cách giữa các số trên mặt đồng hồ ứng với $5$ phút ($5 \times 1 = 5$ phút, $5 \times 6 = 30$ phút).
\end{itemize}
}
\end{tcolorbox}
\end{minipage}

\vspace{-1.5mm}
\begin{tcolorbox}[
  enhanced,
  colback=orange!4!white,
  colframe=orange!80!black,
  arc=2mm, boxrule=0.8pt,
  left=2mm, right=2mm, top=1.5mm, bottom=1.5mm,
  title={\footnotesize\bfseries\color{white}\faSitemap\ 3. Sơ đồ tư duy: Bậc thang dãy tích Bảng nhân 5}
]
\centering
\begin{tikzpicture}[x=0.9cm, y=0.6cm]
  \foreach \i/\val in {1/5, 2/10, 3/15, 4/20, 5/25, 6/30, 7/35, 8/40, 9/45, 10/50} {
    \node[circle, draw=teal!80!black, fill=teal!15, inner sep=2pt, font=\tiny\bfseries] at (\i*1.1, 0) {\val};
    \ifnum\i<10
      \draw[->, thick, orange!80!black] (\i*1.1+0.3, 0.2) to[bend left=40] node[above, font=\tiny] {+5} (\i*1.1+0.8, 0.2);
    \fi
  }
\end{tikzpicture}
\end{tcolorbox}

\vspace{-1.5mm}
\noindent
\begin{minipage}[t]{0.49\textwidth}\vspace{0pt}
\begin{tcolorbox}[
  enhanced,
  colback=green!4!white,
  colframe=green!60!black,
  arc=2mm, boxrule=0.8pt,
  left=2mm, right=2mm, top=1.8mm, bottom=1.8mm,
  title={\footnotesize\bfseries\color{white}\faCheckCircle\ 4. Ví dụ thực hành chuẩn mực}
]
{\footnotesize
\textbf{1. Tính nhẩm:}\\
$5 \times 3 = \mathbf{15}$; $5 \times 6 = \mathbf{30}$; $5 \times 8 = \mathbf{40}$.
\textbf{2. Bài toán lời văn:}\\
Mỗi ngôi sao có $5$ cánh. Hỏi $6$ ngôi sao như thế có tất cả bao nhiêu cánh?\\
\textit{Bài giải:}\\
$6$ ngôi sao có tất cả số cánh là:
$$5 \times 6 = 30\text{ (cánh)}$$
Đáp số: $30$ cánh sao.
}
\end{tcolorbox}
\end{minipage}
\hfill
\begin{minipage}[t]{0.49\textwidth}\vspace{0pt}
\begin{tcolorbox}[
  enhanced,
  colback=red!4!white,
  colframe=red!75!black,
  arc=2mm, boxrule=0.8pt,
  left=2mm, right=2mm, top=1.8mm, bottom=1.8mm,
  title={\footnotesize\bfseries\color{white}\faExclamationTriangle\ 5. Mẹo tránh bẫy trắc nghiệm}
]
{\footnotesize
\begin{itemize}\setlength{\itemsep}{1pt}\setlength{\leftskip}{-2.5mm}
  \item \textbf{Nhận dạng nhanh chữ số tận cùng:}\\
  Nếu nhân với số chẵn ($2, 4, 6, 8, 10$) $\to$ tận cùng là \textbf{0}.\\
  Nếu nhân với số lẻ ($1, 3, 5, 7, 9$) $\to$ tận cùng là \textbf{5}.
  \item Không bao giờ có kết quả tận cùng là $1, 2, 3, 4, 6, 7, 8, 9$.
\end{itemize}
}
\end{tcolorbox}
\end{minipage}


\end{document}
